%-------------------------
% General Academic CV
%------------------------

\documentclass[letterpaper,11pt]{article}

\usepackage[empty]{fullpage}
\usepackage{titlesec}
\usepackage{enumitem}
\usepackage[hidelinks]{hyperref}
\usepackage{fancyhdr}
\usepackage[english]{babel}
\usepackage{tabularx}
\usepackage{microtype}
\usepackage{charter}
\usepackage[T1]{fontenc}
\usepackage{xcolor}

\definecolor{rulegray}{gray}{0.80}
\definecolor{textgray}{gray}{0.12}

\pagestyle{fancy}
\fancyhf{}
\fancyfoot{}
\renewcommand{\headrulewidth}{0pt}
\renewcommand{\footrulewidth}{0pt}
\setlength{\footskip}{12pt}

% Margins
\addtolength{\oddsidemargin}{-0.5in}
\addtolength{\evensidemargin}{0in}
\addtolength{\textwidth}{1in}
\addtolength{\topmargin}{-.55in}
\addtolength{\textheight}{1.1in}

\urlstyle{same}
\raggedbottom
\raggedright
\setlength{\tabcolsep}{0in}

% Clean academic section styling
\titleformat{\section}{
    \large\bfseries\scshape\raggedright
}{}{0em}{}[\color{rulegray}\titlerule\vspace{-3pt}]

\pdfgentounicode=1
\color{textgray}

%-------------------------
% Custom commands
\newcommand{\resumeSubheading}[4]{
  \vspace{-1pt}\item
    \begin{tabular*}{\textwidth}[t]{l@{\extracolsep{\fill}}r}
      \textbf{#1} & {\small #2}\vspace{1pt}\\
      \textit{#3} & {\small #4}\\
    \end{tabular*}\vspace{-4pt}
}

\newcommand{\resumeSubHeadingListStart}{\begin{itemize}[leftmargin=0in, label={}, itemsep=6pt, topsep=3pt]}
\newcommand{\resumeSubHeadingListEnd}{\end{itemize}}
\newcommand{\resumeItemListStart}{\begin{itemize}[leftmargin=0.18in, itemsep=2pt, topsep=2pt]}
\newcommand{\resumeItemListEnd}{\end{itemize}\vspace{-2pt}}
\newcommand{\resumeItem}[1]{\item\small{{#1 \vspace{-1pt}}}}

%-------------------------------------------
\begin{document}

\begin{center}
    {\LARGE \textbf{Egbert Castro, PhD}} \\ \vspace{6pt}
    \small
    \href{https://egbertcastro.com}{egbertcastro.com}
    \hspace{5pt}$\cdot$\hspace{5pt}
    \href{mailto:egbcastro@gmail.com}{egbcastro@gmail.com}
    \hspace{5pt}$\cdot$\hspace{5pt}
    \href{https://www.linkedin.com/in/castroegbert}{linkedin.com/in/castroegbert}
    \hspace{5pt}$\cdot$\hspace{5pt}
    \href{https://github.com/egbert}{github.com/egbert}
    \hspace{5pt}$\cdot$\hspace{5pt}
    Los Angeles, CA
\end{center}

\section{Research Interests}
{\small
Machine learning for biology and scientific discovery; generative modeling for proteins and molecules; representation learning for structured biological data; latent-space optimization and controllable generation; graph neural networks and geometric deep learning.
}

\section{Education}
\resumeSubHeadingListStart
  \resumeSubheading
    {Yale University}{New Haven, CT}
    {PhD, Computational Biology \& Bioinformatics}{2018 -- 2023}
    \resumeItemListStart
      \resumeItem{Advisor: Prof.\ Smita Krishnaswamy}
      \resumeItem{Dissertation: \textit{Representation Learning and Generative Modeling for Biomolecular Sequences and Structures}}
    \resumeItemListEnd

  \resumeSubheading
    {University of California, San Diego}{La Jolla, CA}
    {BS, Pharmacological Chemistry; Minor in Mathematics}{2014 -- 2018}
\resumeSubHeadingListEnd

\section{Teaching Experience}
\resumeSubHeadingListStart
  \resumeSubheading
    {Yale University}{New Haven, CT}
    {Teaching Assistant}{}
    \resumeItemListStart
      \resumeItem{\textbf{CPSC 4520b, Deep Learning Theory and Applications} (Instructor: Dr.\ Smita Krishnaswamy). Supported a course covering the principles and practice of neural networks and modern deep learning, with applications in computer vision, natural language processing, and biomedical problems.}
      \resumeItem{\textbf{MCDB 570, Biotechnology} (Instructor: Dr.\ Craig Crews). Supported instruction on cellular, molecular, and chemical techniques in biotechnology, including applications to therapeutics, industrial agents, and biological research.}
    \resumeItemListEnd
\resumeSubHeadingListEnd

\section{Research and Professional Experience}
\resumeSubHeadingListStart
  \resumeSubheading
    {Ascent Bio}{Hybrid - Los Angeles, CA and New York, NY}
    {Co-founder, AI Research}{2023 -- Present}
    \resumeItemListStart
      \resumeItem{Led the zero-to-one effort to build a VC-backed agentic science platform, from research to production.}
      \resumeItem{Accepted into the Merck Digital Sciences Accelerator and IndieBio NYC Accelerator.}
      \resumeItem{Designed and built agentic systems for scientific and drug-discovery workflows that combined foundation models, scientific tools, and code execution across multi-step research tasks.}
      \resumeItem{Developed evaluation frameworks for correctness, tool-use reliability, and failure-mode analysis on realistic scientific problems.}
      \resumeItem{Contributed to molecule reasoning-model research using synthetic data generation, reproducible experimentation, and rapid iteration in a shared research codebase.}
    \resumeItemListEnd

  \resumeSubheading
    {Yale University -- Krishnaswamy Lab}{New Haven, CT}
    {Graduate Researcher, Computational Biology \& Bioinformatics}{2018 -- 2023}
    \resumeItemListStart
      \resumeItem{Conducted research in geometric deep learning and generative modeling for biological sequences, structures, and graphs, with applications spanning proteins, RNA, molecular graphs, and clinical prediction.}
      \resumeItem{Developed transformer-based protein generative models with regularized latent-space optimization for gradient-guided search in learned sequence representations; first-author work published in \textit{Nature Machine Intelligence}.}
      \resumeItem{Designed representation-learning methods for RNA secondary-structure landscapes and molecular graph generation, integrating graph-based architectures, geometric scattering, and task-specific biological objectives.}
      \resumeItem{Established reproducible multi-GPU training pipelines, curated benchmark datasets, and defined biologically grounded evaluation protocols for comparative analysis of biomolecular machine-learning methods.}
    \resumeItemListEnd

  \resumeSubheading
    {Genentech}{South San Francisco, CA}
    {Machine Learning Research Intern, Discovery Chemistry}{Summer 2018}
    \resumeItemListStart
      \resumeItem{Developed predictive models for ADMET properties on internal compound libraries using deep neural networks.}
      \resumeItem{Built automated machine-learning pipelines integrating RDKit-based featurization with model training and evaluation.}
      \resumeItem{Partnered with discovery scientists on model evaluation for medicinal chemistry use cases.}
    \resumeItemListEnd
\resumeSubHeadingListEnd

\section{Publications}
{\small \textit{$\dagger$ denotes first author.}}
\vspace{4pt}

{\small
\textbf{Peer-Reviewed Journal Articles and Conference Papers}
\begin{itemize}[leftmargin=0.18in, itemsep=5pt, topsep=2pt]
  \item[$\dagger$] \textbf{Castro, E.}, Godavarthi, A., Rubinfien, J., Givechian, K., Bhaskar, D., Krishnaswamy, S.
  \textit{Transformer-based protein generation with regularized latent space optimization}.
  \textbf{Nature Machine Intelligence}, 4(10):840--851, 2022.

  \item Bhaskar, D., Grady, J., \textbf{Castro, E.}, Perlmutter, M., Krishnaswamy, S.
  \textit{Molecular graph generation via geometric scattering}.
  2022 IEEE 32nd International Workshop on Machine Learning for Signal Processing (MLSP), 1--6, 2022.

  \item Shung, D., Huang, J., \textbf{Castro, E.}, Tay, J.K., Simonov, M., Laine, L., Batra, R., Krishnaswamy, S.
  \textit{Neural network predicts need for red blood cell transfusion for patients with acute gastrointestinal bleeding admitted to the intensive care unit}.
  \textbf{Scientific Reports}, 11(1):8827, 2021.

  \item[$\dagger$] \textbf{Castro, E.}, Benz, A., Tong, A., Wolf, G., Krishnaswamy, S.
  \textit{Uncovering the folding landscape of RNA secondary structure using deep graph embeddings}.
  2020 IEEE International Conference on Big Data (Big Data), 4519--4528, 2020.

  \item Tew, B.Y., Hong, T.B., Otto-Duessel, M., Elix, C., \textbf{Castro, E.}, He, M., Wu, X., Pal, S.K., Kalkum, M., Jones, J.O.
  \textit{Vitamin K epoxide reductase regulation of androgen receptor activity}.
  \textbf{Oncotarget}, 8(8):13818--13831, 2017.
\end{itemize}

\vspace{4pt}
\textbf{Patent Applications}
\begin{itemize}[leftmargin=0.18in, itemsep=5pt, topsep=2pt]
  \item[$\dagger$] \textbf{Castro, E.}, Krishnaswamy, S., Godavarthi, A., Rubinfien, J.
  \textit{Regularized deep learning based improvement of biomolecules}.
  US Patent Application 18/061,272, published September 14, 2023.
\end{itemize}

\vspace{4pt}
\textbf{Preprints}
\begin{itemize}[leftmargin=0.18in, itemsep=5pt, topsep=2pt]
  \item Afrasiyabi, A., Kovalic, J., Liu, C., \textbf{Castro, E.}, Weinreb, A., Varol, E., Miller III, D.M., Hammarlund, M., Krishnaswamy, S.
  \textit{CellSpliceNet: Interpretable multimodal modeling of alternative splicing across neurons in C.\ elegans}.
  \textbf{bioRxiv}, 2025.06.22.660966, 2025.

  \item Viswanath, S., Bhaskar, D., Johnson, D.R., Rocha, J.F., \textbf{Castro, E.}, Grady, J.D., Grigas, A.T., Perlmutter, M.A., O'Hern, C.S., Krishnaswamy, S.
  \textit{ProtSCAPE: Mapping the landscape of protein conformations in molecular dynamics}.
  arXiv preprint arXiv:2410.20317, 2024.
\end{itemize}

\vspace{4pt}
\textbf{Conference Abstracts}
\begin{itemize}[leftmargin=0.18in, itemsep=5pt, topsep=2pt]
  \item[$\dagger$] \textbf{Castro, E.}, Shung, D.
  \textit{993 NEURAL NETWORK PREDICTS DROP IN HEMOGLOBIN REQUIRING TRANSFUSION FOR PATIENTS WITH ACUTE GASTROINTESTINAL BLEEDING ADMITTED TO THE ICU}.
  \textbf{Gastroenterology}, 158(6):S-196--S-197, 2020.
\end{itemize}
}

\section{Presentations \& Talks}
{\small
\begin{itemize}[leftmargin=0.18in, itemsep=4pt, topsep=2pt]
  \item \textbf{ProtSCAPE: Mapping the landscape of protein conformations in molecular dynamics}. SMT Conference, 2025.
  \item \textbf{ProGSNN}. Presented at the Learning Meaningful Representations of Life Workshop at ICLR, 2021, and at the NLM Training Grant Conference, 2022.
  \item \textbf{Molecular graph generation via geometric scattering}. IEEE International Workshop on Machine Learning for Signal Processing, 2022.
  \item \textbf{Guided Generative Protein Design using Regularized Transformers}. ISMB/ECCB, 2021.
  \item \textbf{Uncovering the Folding Landscape of RNA Secondary Structure with Deep Graph Embeddings}. Presented at the GRLB Workshop at ICML and at the IEEE International Conference on Big Data, 2020.
\end{itemize}
}

\section{Technical Skills}
\vspace{2pt}
{\small
\begin{tabularx}{\textwidth}{@{}>{\bfseries}l @{\hspace{12pt}} X@{}}
Methods & transformers, generative models, latent-space optimization, representation learning, sequence modeling, graph neural networks, geometric deep learning \\
Domains & protein sequence modeling, molecular property prediction, RNA structure, multimodal biological data, AI for scientific discovery \\
Tools & Python, PyTorch, JAX, Hugging Face, scikit-learn, RDKit, Weights and Biases, multi-GPU training, benchmarking, reproducible experimentation \\
\end{tabularx}
}

\end{document}
